\documentclass[10pt,a4paper]{article}
\usepackage[utf8]{inputenc}
\usepackage[T1]{fontenc}
\usepackage[a4paper,top=20mm,bottom=20mm,left=16mm,right=16mm,headheight=14pt]{geometry}
\usepackage{fancyhdr}
\usepackage{booktabs}
\usepackage{longtable}
\usepackage{array}
\usepackage{enumitem}
\usepackage{amssymb}
\usepackage{xcolor}
\usepackage[hyphens]{url}
\usepackage[colorlinks=true,urlcolor=blue!60!black,linkcolor=blue!60!black]{hyperref}

\setlength{\parindent}{0pt}
\setlength{\parskip}{3pt}
\renewcommand{\arraystretch}{1.28}

\pagestyle{fancy}
\fancyhf{}
\fancyhead[L]{\small LANGUAGE \& SOCIETY}
\fancyhead[R]{\small TRANSCRIPT ANALYSIS WORKSHEET}
\fancyfoot[C]{\small\thepage}
\renewcommand{\headrulewidth}{0.4pt}

\begin{document}

\begin{center}
    {\large\bfseries LANGUAGE \& SOCIETY PROJECT}\\[3pt]
    {\LARGE\bfseries Transcript Analysis Worksheet}\\[3pt]
    {\normalsize Qualitative analysis of political news discourse}
\end{center}
\vspace{2pt}
\noindent\rule{\textwidth}{0.4pt}
\vspace{4pt}


\section*{Creator 6: The Wire (Independent Left-Liberal Long-Form Digital Journalism)}
\begin{itemize}[leftmargin=1.5em, itemsep=3pt, topsep=4pt]
  \item \textbf{Analyst:} Baani Solanki (\texttt{2025114016})
  \item \textbf{Channel / Outlet:} The Wire (\texttt{@TheWire}, Foundation for Independent Journalism)
  \item \textbf{Ideological / Editorial Positioning:} Independent, critical left-liberal English digital journalism; uses long-form one-on-one elite interviews (with senior historians/authors and cultural figures) to situate breaking controversies within structural critiques of Hindutva/Sangh Parivar institutional history, PMO centralization, and democratic backsliding.
  \item \textbf{Transcript Files Referenced:}
  \begin{itemize}[leftmargin=1.4em, itemsep=2pt, topsep=2pt]
    \item Event 1: \texttt{transcripts/event1\_ram\_mandir/06\_the\_wire\_7rUL8fGgRJA.txt}
    \item Event 2: \texttt{transcripts/event2\_neet\_protest/06\_the\_wire\_pYaTVvHeGoM.txt}
  \end{itemize}
\end{itemize}
\vspace{4pt}\noindent\rule{\textwidth}{0.4pt}\vspace{4pt}

\section*{Guidelines: Five features to look for in every event}
\begin{longtable}{@{}>{\raggedright\arraybackslash}p{0.17\textwidth} >{\raggedright\arraybackslash}p{0.20\textwidth} >{\raggedright\arraybackslash}p{0.24\textwidth} >{\raggedright\arraybackslash}p{0.32\textwidth}@{}}
\toprule
\textbf{Feature} & \textbf{What to look for} & \textbf{Illustrative example} & \textbf{Questions to answer} \\
\midrule
\endhead
\textbf{1. Naming} & What are people, groups, actions, or events called? & "Students", "protesters", "agitators"; "irregularity", "scam", "controversy". & What exact label is used? Who or what does it refer to? Who uses it? What alternative label could have been used? \\
\addlinespace[4pt]
\textbf{2. Presupposition} & What is treated as already true or taken for granted? & "Why are students creating chaos again?" may presuppose earlier disruption. & What does the wording assume? Is that assumption supported, disputed, or left unexamined? Does it appear in a question or words such as "again", "still", or "finally"? \\
\addlinespace[4pt]
\textbf{3. Agency} & Who is shown doing something? Who is acted upon? Who disappears? & "Police injured protesters" vs. "Protesters were injured." & Who is the actor? Who experiences the action? Is responsibility named, backgrounded, or omitted? \\
\addlinespace[4pt]
\textbf{4. Antagonisation} & Who is positioned against whom? & "Ordinary citizens" vs. "political troublemakers". & Which groups are contrasted? Is one presented as more legitimate, moral, patriotic, dangerous, or disloyal? Is an "us vs. them" division implied? \\
\addlinespace[4pt]
\textbf{5. Intensification / softening} & What is exaggerated, minimised, softened, or euphemised? & "Violent mob" vs. "peaceful demonstrators"; "incident" instead of a direct description. & Which words heighten emotion or seriousness? Which minimise it? Are metaphors, dramatic adjectives, hedges, or euphemisms used? Does the wording match the evidence? \\
\bottomrule
\end{longtable}
\textit{Reminder: A label or implication is not automatically false or manipulative. Examine the surrounding context and evidence.}

\vspace{4pt}\noindent\rule{\textwidth}{0.4pt}\vspace{4pt}

\section*{Event1. Ram Mandir donation controversy}
\textbf{Focus:} How is the controversy framed---as financial accountability, religious sentiment, political motivation, or something else? Be precise about the specific allegation or episode covered.

\begin{itemize}[leftmargin=1.5em, itemsep=3pt, topsep=4pt]
  \item \textbf{Outlet:} The Wire (Special Discussion hosted by M.K. Venu / Venugopal with senior journalist \& Ayodhya historian Nilanjan Mukhopadhyay)
  \item \textbf{Date:} \texttt{2026-07-01}
  \item \textbf{Video title:} \textit{Ram Temple Donations Loot: Biggest Credibility Test For Modi-RSS} (\url{https://youtu.be/7rUL8fGgRJA?si=9HV3jAYEipsTwvAI})
  \item \textbf{Timestamp(s):} Relevant Event 1 segment \texttt{[00:00--39:06]} (\texttt{[00:00--02:48]} cold-open highlights reel + \texttt{[03:00--39:06]} full conversation; \texttt{100\%} focused on Event 1; total duration \texttt{39:18}).
  \item \textbf{Analyst:} Baani Solanki (\texttt{2025114016})
  \item \textbf{Language:} English (Native English long-form discussion, Roman script; non-AI-dubbed)
\end{itemize}

\subsection*{Event-specific questions (Event 1)}
\begin{enumerate}[leftmargin=1.6em, itemsep=4pt, topsep=4pt]
  \item \textbf{What terms are used: "scam", "alleged irregularity", "corruption", "political controversy", or something else?}
  \begin{itemize}[leftmargin=1.4em, itemsep=2pt, topsep=2pt]
    \item Uses high-level structural, legal, and moral terminology: \texttt{"Ram Temple Donations Loot"} (\texttt{[Title]}), \texttt{"the Ram Temple scandal"} (\texttt{[03:10]}), \texttt{"massive theft or decoity [dacoity], as Nripendra Misra... has described it"} (\texttt{[00:09--00:17]}, \texttt{[04:23--04:32]}), \texttt{"a land scam of this magnitude involving potentially massive money laundering via land"} (\texttt{[00:17--00:26]}, \texttt{[04:32--04:40]}), \texttt{"hawala"} (\texttt{[26:42]}), \texttt{"moral degeneration"} (\texttt{[00:44]}, \texttt{[09:08]}), \texttt{"moral and ethical morass"} (\texttt{[28:12--28:21]}), and \texttt{"a platform of corruption, theft and lies"} (\texttt{[37:48]}).
  \end{itemize}
  \item \textbf{How are the temple, trust, donors, and people making allegations described?}
  \begin{itemize}[leftmargin=1.4em, itemsep=2pt, topsep=2pt]
    \item \textbf{Temple \& Movement:} Analysed historically not as an eternal pilgrimage centre, but as \texttt{"the Sangha Parivar's, the BJP's... ideological mission, project"} since \texttt{1949}/\texttt{1983} (\texttt{[03:37--03:48]}, \texttt{[05:29--05:38]}) and \texttt{"Ayodhya the Hindu Vatican... actually the Sangh Parivar's creation"} (\texttt{[35:04--35:13]}).
    \item \textbf{Trust \& Leadership:} Described as a government-nominated \texttt{"corporate bureaucratic structure"} (\texttt{[11:00--11:39]}, \texttt{[23:19]}) staffed by \texttt{"political loyalists"} (\texttt{[25:39--25:48]}) from an RSS culture of \texttt{"shadow organizations"} (\texttt{[22:24]}) and unaccounted \textit{Guru Dakshina} (\texttt{[17:35--18:40]}). Champat Rai is characterized via his nickname \texttt{"Patwari"} (adept at land/revenue records rather than transparent corporate governance, \texttt{[24:18--24:36]}), while Construction Chair Nripendra Misra is identified as \texttt{"very much a part of the core team of Mr. Modi"} who \texttt{"surely knew"} and kept \texttt{"superiors [PMO] in the loop"} (\texttt{[11:48--14:55]}).
    \item \textbf{Donors:} Portrayed with sociological nuance---on one hand sincere devotees from the 1989 \textit{Ram Shila Yatras} onward who \texttt{"put their body and soul into this movement"} only to find \texttt{"the darshan has been hollowed out, that it has been stolen"} (\texttt{[36:28--37:08]}), and on the other hand rational-losing \texttt{"bhakts"} viewing Modi as \texttt{"not just a biological being"} who might \texttt{"overlook anything"} (\texttt{[02:18--02:27]}, \texttt{[35:43--36:20]}).
    \item \textbf{People Making Allegations:} Highlights \textbf{internal Hindu/movement critics}---specifically dissenting orders like \texttt{Nirmohi Akhada} (\texttt{[ASR: Nirmala Akada]}, \texttt{[29:11--29:38]}) calling Sangh Parivar trustees \texttt{"all thieves"} after being excluded from the Trust, alongside historical references to honest 1990 Income Tax Deputy Commissioner \texttt{Vishwa Bandhu Gupta} (\texttt{[ASR: Vishva Hindu Gupta]}, \texttt{[15:33--17:15]}) and Karnataka minister \texttt{Priyank Kharge} (\texttt{[18:59--19:09]}).
  \end{itemize}
  \item \textbf{Is the issue framed as financial accountability, an attack on Hindu faith, a political dispute, or a matter of public trust?}
  \begin{itemize}[leftmargin=1.4em, itemsep=2pt, topsep=2pt]
    \item Framed as an \textbf{existential ideological and moral credibility crisis (\texttt{"Biggest Credibility Test For Modi-RSS"}, \texttt{"biggest challenge for the RSS after the assassination of Mahatma Gandhi"}, \texttt{[08:23--08:31]})} rooted in the Sangh Parivar's 40-year institutional refusal to maintain transparent, audited financial accounts (\texttt{[17:24--22:32]}).
  \end{itemize}
  \item \textbf{Are questions about donations presented as legitimate scrutiny or as an attempt to malign the temple or its supporters?}
  \begin{itemize}[leftmargin=1.4em, itemsep=2pt, topsep=2pt]
    \item \textbf{Unequivocally legitimate structural and forensic scrutiny}---calling for Enforcement Directorate (\texttt{ED}) investigation into land-sale \textit{hawala} trails (\texttt{[26:42--27:18]}) and auditing of unregistered Sangh affiliates (\texttt{[17:24--21:30]}).
  \end{itemize}
  \item \textbf{Are allegations attributed to specific people and supported with evidence, or presented as established facts?}
  \begin{itemize}[leftmargin=1.4em, itemsep=2pt, topsep=2pt]
    \item \textbf{Carefully hedged on financial destinations, deductive on PMO awareness, and richly documented on institutional history:} Host M.K. Venu attributes \texttt{"massive theft or dacoity"} directly to Nripendra Misra (\texttt{[04:23]}), hedges money-laundering beneficiaries with epistemic modals (\texttt{"potentially massive money laundering... We don't know who are the beneficiaries... They could be political coffers, they could be Sangh Parivar associates"}, \texttt{[04:40--04:59]}), and explicitly notes \texttt{"We do not have any empirical data as yet"} on voter outrage (\texttt{[37:08--37:19]}). Meanwhile, Nilanjan Mukhopadhyay uses \textbf{deductive bureaucratic reasoning} (\texttt{"I find it very hard to believe... on that I have absolutely no doubt"}, \texttt{[13:04--14:55]}) to infer PMO knowledge via Nripendra Misra, backed by deep archival knowledge of \texttt{1983--1990} VHP/\texttt{Mukti Yagna Samiti}/\texttt{Ram Janmabhoomi Nyas} structures (\texttt{[15:25--22:50]}).
  \end{itemize}
  \item \textbf{How are denials and counterarguments represented: explained, dismissed, challenged, or omitted?}
  \begin{itemize}[leftmargin=1.4em, itemsep=2pt, topsep=2pt]
    \item \textbf{Analysed politically rather than hosted live} (as there is no BJP/RSS guest in this two-journalist format): UP CM Yogi Adityanath's public assurance is dismissed as \texttt{"having to put up a brave front"} (\texttt{[04:14]}, \texttt{[29:00]}) amid \texttt{"a major fissure between Yogi Adityanath and the central leadership"} (\texttt{[30:04--30:24]}), and the RSS's traditional defense---playing state-persecuted victimhood as during the Emergency or 1992 ban---is dismantled because \texttt{"Today, they are in power. So, they can't play victims"} (\texttt{[10:14--10:49]}).
  \end{itemize}
  \item \textbf{Does the presenter distinguish criticism of the trust or its financial practices from criticism of the temple or religion itself?}
  \begin{itemize}[leftmargin=1.4em, itemsep=2pt, topsep=2pt]
    \item \textbf{Distinguishes classical Hinduism from Sangh Parivar political Hindutva:} Mukhopadhyay emphasizes that established Hindu spiritual heads (\texttt{Shankaracharyas} / traditional \textit{acharyas} and \textit{akhadas} in Varanasi, Mathura, Haridwar) were deliberately \texttt{"kept out"} by the Sangh Parivar when constructing \texttt{"Ayodhya the Hindu Vatican"} (\texttt{[34:07--35:22]}), framing the scandal as a crisis of \textit{Sangh Parivar real-estate/corporate Hindutva} (\texttt{[26:14--26:32]}), not traditional Hindu spirituality.
  \end{itemize}
  \item \textbf{Who is portrayed as protecting donors' interests, and who as exploiting religious sentiment?}
  \begin{itemize}[leftmargin=1.4em, itemsep=2pt, topsep=2pt]
    \item \textbf{Exploiting religious sentiment:} The BJP-RSS-VHP \textit{"Sangh Parivar"}, which mobilized grassroots \textit{Ram Shila} devotees (\texttt{[36:28--37:08]}), sidelined traditional Hindu \textit{akhadas} (\texttt{[29:20--29:38]}, \texttt{[34:07--35:22]}), and deliberately avoided \texttt{"leak-proof"} accounting structures (\texttt{[24:44--25:01]}) while getting \texttt{"addicted to power"} and plush Delhi headquarters (\texttt{[31:22--32:27]}).
    \item \textbf{Protecting public accountability:} Transferred 1990 tax officer Vishwa Bandhu Gupta (\texttt{[15:33--17:15]}) and independent journalists/opposition actors if they can \texttt{"think creatively"} ahead of UP elections (\texttt{[37:28--38:24]}).
  \end{itemize}
\end{enumerate}

\subsection*{Apply the five shared features (Event 1)}
\begin{longtable}{@{}>{\raggedright\arraybackslash}p{0.19\textwidth} >{\raggedright\arraybackslash}p{0.38\textwidth} >{\raggedright\arraybackslash}p{0.37\textwidth}@{}}
\toprule
\textbf{Feature} & \textbf{Exact wording + timestamp} & \textbf{What it does in context} \\
\midrule
\endhead
\textbf{1. Naming} & \texttt{"shadow banks, you know, shadow organizations"} \texttt{[22:24]}; \texttt{"corporate bureaucratic structure"} staffed by \texttt{"political loyalists"} \texttt{[23:19]}, \texttt{[25:48]}; \texttt{"Ayodhya the Hindu Vatican"} \texttt{[35:13]} & Replaces devotional labels with critical political-economy and ecclesiastical terms (\texttt{shadow organizations}, \texttt{corporate bureaucratic structure}, \texttt{Hindu Vatican}), demystifying the Trust as a centralized Sangh power apparatus. \\
\addlinespace[4pt]
\textbf{2. Presupposition} & \texttt{"within 2 years of that, the unraveling of the Ram temple project in moral terms, in ethical terms"} \texttt{[06:15--06:26]} \newline  \texttt{"Sangha Parivar values got obviously got completely contaminated by the real estate game"} \texttt{[26:14]} & Nominalization (\texttt{the unraveling of the Ram temple project}) and adverb \texttt{"obviously"} presuppose that the temple project has already morally collapsed and been corrupted by real-estate profiteering. \\
\addlinespace[4pt]
\textbf{3. Agency} & \texttt{"Setting up structures was not done. It was probably not done because... they wanted that it should not be a leak-proof system. So you keep a certain amount of ambiguity, a certain amount of fluidity."} \texttt{[24:36--25:01]} & Converts what other outlets call passive "oversight" or "informality" (\texttt{anaupcharik} on Aaj Tak) into \textbf{deliberate, strategic agency} (\texttt{they wanted that it should not be a leak-proof system}). \\
\addlinespace[4pt]
\textbf{4. Antagonisation} & \texttt{"RSS chief described it as the real freedom India has attained as opposed to what was attained in 1947... and within two years... the entire edifice was built on a platform of corruption, theft and lies."} \texttt{[05:57--06:15]}, \texttt{[37:48]} & Contrasts RSS Chief Mohan Bhagwat's grand civilizational claim (\texttt{real freedom} vs \texttt{1947}) against sordid financial crime (\texttt{hawala}, \texttt{corruption, theft and lies}), framing the Sangh Parivar's moral claim over India as hollow. \\
\addlinespace[4pt]
\textbf{5. Intensification / softening} & \textbf{Intensification:} \texttt{"surrounded by 100 questions today, maybe 1,000 questions"} \texttt{[00:00]}, \texttt{[04:05]}; \texttt{"biggest challenge for the RSS after the assassination of Mahatma Gandhi"} \texttt{[08:23--08:31]} \newline  \textbf{Epistemic Softening:} \texttt{"potentially massive money laundering... We don't know who are the beneficiaries... They could be..."} \texttt{[04:40--04:59]}; \texttt{"We do not have any empirical data as yet"} \texttt{[37:19]} & Pairs sweeping historical intensification (comparing the crisis to 1948 post-Gandhi assassination) with careful journalistic hedging (\texttt{potentially}, \texttt{could be}, \texttt{no empirical data}) wherever hard forensic proof is still missing. \\
\bottomrule
\end{longtable}

\subsection*{Short summary (Event 1)}
\begin{itemize}[leftmargin=1.5em, itemsep=3pt, topsep=4pt]
  \item \textbf{What interpretation of the event does this video encourage?} That the Ayodhya donation \texttt{"dacoity"} and \texttt{5x--10x} land-price inflation (\texttt{[04:23]}, \texttt{[26:42]}) are not rogue clerk anomalies, but the direct institutional consequence of the Sangh Parivar's decades-long ethos of opaque, unaudited \texttt{"shadow organizations"} (\texttt{[17:24--22:32]}, \texttt{[24:44--25:01]}) colliding with massive post-2020 state patronage under PMO-linked appointees like Nripendra Misra and \texttt{"Patwari"} Champat Rai (\texttt{[11:48--14:55]}, \texttt{[24:18]}).
  \item \textbf{Whose perspective is foregrounded, challenged, or missing?} \textit{Foregrounded:} Secular-historical political analysis (M.K. Venu \& Nilanjan Mukhopadhyay), sidelined Hindu monastic orders (\texttt{Nirmohi Akhada}, traditional spiritual heads at \texttt{[29:11--29:38]}, \texttt{[34:07--35:22]}), and ordinary 1989 \textit{Ram Shila} devotees whose darshan was \texttt{"hollowed out"} (\texttt{[36:28--37:08]}). \textit{Challenged:} Sangh Parivar moral supremacy (\texttt{[05:57--06:35]}), RSS financial opacity (\texttt{[17:24--19:35]}), and PMO/Yogi plausible deniability (\texttt{[12:08--14:55]}, \texttt{[28:31--30:24]}). \textit{Missing:} Any Trust/RSS/BJP representative or ground reportage from Ayodhya's counting room.
  \item \textbf{What evidence supports your interpretation? What complicates it?} \textit{Supports:} Documented historical precedents---\texttt{1990} IT notice by Vishwa Bandhu Gupta over \texttt{Rs.\,700 crore} VHP collections and his punitive transfer (\texttt{[15:33--17:15]}), dual registered/unregistered 1984 \texttt{Mukti Yagna Samiti} vs \texttt{Nyas} architecture (\texttt{[20:36--22:24]}), \texttt{2019} Supreme Court trust mandate (\texttt{[11:08--11:39]}), and Nripendra Misra's own word \texttt{"dacoity"} (\texttt{[04:23]}). \textit{Complicates:} Both speakers candidly admit they lack direct proof of money-laundering recipients (\texttt{[04:48]}), ED investigation trails (\texttt{[26:51--27:18]}), or empirical data on whether \texttt{"bhakts"} will actually punish the BJP electorally in UP (\texttt{[36:01--37:28]}).
\end{itemize}
\vspace{4pt}\noindent\rule{\textwidth}{0.4pt}\vspace{4pt}

\section*{Event2. Students' protest regarding NEET paper leak}
\textbf{Focus:} How does coverage construct the legitimacy of student protest and responsibility for the underlying problem?

\begin{itemize}[leftmargin=1.5em, itemsep=3pt, topsep=4pt]
  \item \textbf{Outlet:} The Wire (\textit{The Interview with Karan Thapar} --- Guest: veteran actor \& educator Naseeruddin Shah)
  \item \textbf{Date:} \texttt{2026-07-22} \textit{(two days after the 20 July "Black Monday" crackdown; three days before Dharmendra Pradhan's resignation)}
  \item \textbf{Video title:} \textit{I stand with the protesting students; this is not the India I love; I don't take Modi seriously} (\url{https://youtu.be/pYaTVvHeGoM?si=TPLfYLT3eB6UHh5o})
  \item \textbf{Timestamp(s):} Relevant Event 2 segments \texttt{[00:00--02:45]} (cold-open preview) and \texttt{[03:59--31:05]} (main interview; total duration \texttt{32:26}; viewer-donation appeals at \texttt{[02:46--03:58]} and \texttt{[31:06--32:26]} excluded per project guidelines; note \texttt{[04:25--06:15]} has sparse ASR during playback of Shah's 1m50s video clip).
  \item \textbf{Analyst:} Baani Solanki (\texttt{2025114016})
  \item \textbf{Language:} English with occasional Hindi/Urdu lexical items (\texttt{dandas}, \texttt{bachchas}, \texttt{Pradhan Sevak}, \texttt{Chowkidar}) (Native English audio, Roman script; non-AI-dubbed)
\end{itemize}

\subsection*{Event-specific questions (Event 2)}
\begin{enumerate}[leftmargin=1.6em, itemsep=4pt, topsep=4pt]
  \item \textbf{Are participants called "students", "candidates", "protesters", "agitators", "mobs", or "political pawns"?}
  \begin{itemize}[leftmargin=1.4em, itemsep=2pt, topsep=2pt]
    \item Warmly humanized and idealized throughout by both Karan Thapar and Naseeruddin Shah as \texttt{"protesting students"} (\texttt{[Title]}, \texttt{[04:09]}), \texttt{"people in their teens or their early 20s"} (\texttt{[06:33--06:41]}), \texttt{"our future"} / \texttt{"the children of our future"} (\texttt{[07:42]}, \texttt{[30:19]}), \texttt{"young children"} / \texttt{"these kids"} (\texttt{[08:38]}, \texttt{[09:46]}, \texttt{[20:15]}, \texttt{[25:33]}), and affectionately \texttt{"these bachas [bachchas]"} (\texttt{"These bachas know what they're doing... these bachas are much smarter than us"}, \texttt{[23:18--23:28]}).
  \end{itemize}
  \item \textbf{Are their stated demands explained, summarised, mocked, questioned, or left out?}
  \begin{itemize}[leftmargin=1.4em, itemsep=2pt, topsep=2pt]
    \item \textbf{Clearly explained and morally endorsed, with a nuanced tactical disagreement on Pradhan's resignation:} Karan Thapar summarises their demands---speaking to MPs about exam leaks and the education system (\texttt{[06:23]}), and Education Minister Dharmendra Pradhan's resignation accepting moral responsibility after multiple leaks and suicides (\texttt{[11:55--12:14]}, \texttt{[22:31--22:40]}). Interestingly, while Shah agrees Pradhan \texttt{"should have done so on his own initiative"} (\texttt{[12:14]}), he calls demanding it \texttt{"pointless"} because \texttt{"no way is he going to resign"} and urges pressing systemic demands more urgent than \texttt{"having this idiot to resign"} (\texttt{[12:23--12:50]}) \textit{(a prediction dramatically overturned 3 days later on 25 July!)}.
  \end{itemize}
  \item \textbf{Does the presenter focus on the original grievance, disruption, political affiliations, or public inconvenience?}
  \begin{itemize}[leftmargin=1.4em, itemsep=2pt, topsep=2pt]
    \item Focuses on \textbf{(a) the moral/democratic betrayal of state violence against peaceful youth (\texttt{[06:23--11:55]}), (b) broader syllabus degradation \& ministerial accountability (\texttt{[12:04--14:57]}, \texttt{[19:35--23:37]}), and (c) a scathing character indictment of PM Narendra Modi and authoritarian fear (\texttt{[00:00--02:45]}, \texttt{[15:07--19:01]}, \texttt{[23:47--30:40]})}---explicitly rejecting any disruption frame (\texttt{"It's not as if they were burning effigies or kicking things or destroying stuff"}, \texttt{[25:42--25:51]}).
  \end{itemize}
  \item \textbf{Who is represented as responsible: examination authorities, government, students, political parties, or others?}
  \begin{itemize}[leftmargin=1.4em, itemsep=2pt, topsep=2pt]
    \item \textbf{Prime Minister Narendra Modi} (\texttt{"arrogant autocrats"}, \texttt{[23:47]}; \texttt{"cowardly prime minister"}, \texttt{[26:02]}), \textbf{Education Minister Dharmendra Pradhan} (\texttt{"this idiot"}, Orwellian \texttt{"classic double speak"}, \texttt{[12:42--13:51]}), \textbf{Health Minister J.P. Nadda} (for theatrical hospital photo-ops instead of meeting hunger-striker Sonam Wangchuk, \texttt{[21:03--21:44]}), \textbf{Delhi Police / RAF / privately hired unbadged \texttt{"lumpens"}} (\texttt{[07:15--09:28]}, \texttt{[11:26--11:37]}), and \textbf{silent Indian elites/celebrities} stifling their conscience for state rewards (\texttt{"a dog with a bone in its mouth can't bark"}, \texttt{[29:31--30:00]}).
  \end{itemize}
  \item \textbf{Are questions framed around whether students have legitimate grievances, or why they are causing disruption?}
  \begin{itemize}[leftmargin=1.4em, itemsep=2pt, topsep=2pt]
    \item Karan Thapar's interview questions are \textbf{heavily presuppositional leading questions affirming student legitimacy and state guilt}: e.g., \texttt{"By its response, has the government embarrassed and betrayed India's cherished democracy?"} (\texttt{[06:41--06:49]}), \texttt{"Doesn't that suggest that our political value system has degraded?"} (\texttt{[22:40--22:48]}), and \texttt{"Do peaceful nonviolent Indian protesters protesting in a just cause have the right to walk to parliament...?"} (\texttt{[25:02--25:23]}).
  \end{itemize}
  \item \textbf{Are claims of violence, disorder, or political manipulation attributed to sources and substantiated?}
  \begin{itemize}[leftmargin=1.4em, itemsep=2pt, topsep=2pt]
    \item Karan Thapar anchors key factual claims to \textbf{named mainstream print sources and legal experts}: \textit{Hindustan Times} on hundreds injured (\texttt{[06:33]}), \textit{The Hindu} on pellet guns and shock batons injuring 25-year-old \texttt{Sheikh Irshad Mansuri} with pellets lodged near a neck blood vessel requiring multiple surgeries (\texttt{[10:07--10:38]}), and Senior Advocate \texttt{Kapil Sibal}'s press conference warning that unnamed FIRs could escalate to draconian \texttt{UAPA} / \texttt{BNS Section 195} charges as in the \texttt{CAA} protests (\texttt{[26:41--27:37]}). Conversely, Shah's claims that plainclothes attackers are \texttt{"uneducated... lumpens who've been recruited by being paid"} (\texttt{[11:26--11:37]}) or that PM Modi \texttt{"left by the back gate"} (\texttt{[26:02--26:13]}) rely on viral video impressions and hearsay (\texttt{"witnesses have said"}, \texttt{[21:23]}).
  \end{itemize}
  \item \textbf{How are police and government actions described compared with protesters' actions? Are similar actions evaluated differently?}
  \begin{itemize}[leftmargin=1.4em, itemsep=2pt, topsep=2pt]
    \item Police/state attackers are described as \texttt{"sadistically beating up these students"} (\texttt{[07:15]}), \texttt{"let loose an army of goons on our future"} (\texttt{[07:42]}), \texttt{"splitting young children's heads open, women being attacked and molested"} (\texttt{[08:38--08:47]}), and firing guns reminiscent of \texttt{"Jallianwala Bagh"} (\texttt{[10:48--11:08]}). When discussing \textbf{police casualties}, Shah openly mocks official injury figures (\texttt{"supposedly more policemen injured than students... astounding... something like 600 policemen and about 60 students"}, \texttt{[06:49--07:07]}) and justifies student pushback via a cornered-victim proverb: \texttt{"the students retaliating. It's not surprising. I mean, even a rat fights back when it's pushed against the wall. And these kids are justified"} (\texttt{[09:36--09:46]}).
  \end{itemize}
  \item \textbf{Are students allowed to explain their demands in their own words, or does the presenter mainly speak for them?}
  \begin{itemize}[leftmargin=1.4em, itemsep=2pt, topsep=2pt]
    \item \textbf{Elite proxy advocacy rather than direct student voices:} Because the format is a studio Zoom interview between veteran journalist Karan Thapar and 75-year-old actor Naseeruddin Shah, \textbf{no student speaks directly}; instead, two prominent public intellectuals speak \textit{in solidarity with} (\texttt{"I stand with them"}, \texttt{[20:25]}; \texttt{"We owe it to the children"}, \texttt{[30:19]}) and interpret the youth movement for an English-speaking civic audience.
  \end{itemize}
  \item \textbf{Does the outlet's framing differ between coverage of the NEET paper-leak issue before the protest and coverage once students begin protesting?}
  \begin{itemize}[leftmargin=1.4em, itemsep=2pt, topsep=2pt]
    \item Rather than delegitimizing students once they march to Parliament, The Wire treats the \textbf{20 July march and students' courageous return to Jantar Mantar in bandages on 21 July (\texttt{[19:09--19:18]})} as a historic moral watershed (\texttt{"bigger than anything that's happened before"}, \texttt{[30:09--30:19]}) exposing 12 years of autocratic governance (\texttt{[00:00--00:08]}, \texttt{[23:47--24:07]}).
  \end{itemize}
\end{enumerate}

\subsection*{Apply the five shared features (Event 2)}
\begin{longtable}{@{}>{\raggedright\arraybackslash}p{0.19\textwidth} >{\raggedright\arraybackslash}p{0.38\textwidth} >{\raggedright\arraybackslash}p{0.37\textwidth}@{}}
\toprule
\textbf{Feature} & \textbf{Exact wording + timestamp} & \textbf{What it does in context} \\
\midrule
\endhead
\textbf{1. Naming} & Attackers: \texttt{"an army of goons"} \texttt{[07:42]}, \texttt{"uneducated... lumpens"} \texttt{[11:26]} \newline  Leaders: \texttt{"this idiot"} \texttt{[12:42]}, \texttt{"arrogant autocrats"} \texttt{[23:47]}, \texttt{"rabblerouser"} \texttt{[24:43]}, \texttt{"cowardly prime minister"} \texttt{[26:02]} \newline  Students: \texttt{"our future"}, \texttt{"these kids"}, \texttt{"these bachas"} \texttt{[07:42]}, \texttt{[23:18]} & Inverts Republic TV's naming matrix: \texttt{"goons"} and \texttt{"lumpens"} are applied to unbadged state-aligned attackers, harsh moral-intellectual epithets target ministers/PM, and familial Hindi endearment (\texttt{bachas}) humanizes protesters. \\
\addlinespace[4pt]
\textbf{2. Presupposition} & Karan Thapar \texttt{[13:10]}: \texttt{"Accountable is the one thing the minister refuses to be"} \newline  \texttt{[22:40]}: \texttt{"Doesn't that suggest that our political value system has degraded?"} \newline  \texttt{[23:47]}: \texttt{"both he and his ministers are behaving like arrogant autocrats"} & Negative-interrogative (\texttt{Doesn't that suggest...}) and categorical declaratives embedded inside interview "questions" presuppose ministerial guilt and autocracy as shared consensus between interviewer and guest. \\
\addlinespace[4pt]
\textbf{3. Agency} & \texttt{"to let loose an army of goons on our future"} \texttt{[07:42]}; \texttt{"wrapping him [Wangchuk] up in bed sheets and dragging him away"} \texttt{[21:14--21:23]} vs \texttt{"the students retaliating... even a rat fights back when it's pushed against the wall"} \texttt{[09:36--09:46]} & Assigns sadistic, unprovoked initiatory agency to the state (\texttt{let loose}, \texttt{sadistically beating}, \texttt{dragging him away}), while framing student force strictly as cornered self-defense (\texttt{retaliating}, \texttt{pushed against the wall}). \\
\addlinespace[4pt]
\textbf{4. Antagonisation} & \texttt{"we are living in 1984... 2 minutes of hate has been transformed into 24 hours of hate"} \texttt{[14:02--14:11]} \newline  \texttt{"That house of parliament is ours. It belongs to the kids... This is not the India I was taught to love."} \texttt{[25:33--26:32]} & Pits an Orwellian, anti-intellectual, autocratic regime (\texttt{1984}, \texttt{24 hours of hate}, \texttt{15th century} syllabus at \texttt{[19:53]}) and complicit silent elites (\texttt{dog with a bone}, \texttt{[29:31]}) against courageous, smart youth reclaiming citizens' Parliament. \\
\addlinespace[4pt]
\textbf{5. Intensification / softening} & \textbf{Intensification:} \texttt{"Lie upon lie upon lie upon lie"} \texttt{[01:05]}, \texttt{[16:34]}; \texttt{"broke my heart and made my blood boil"} \texttt{[08:19]}; \texttt{"deaf, dumb and blind"} \texttt{[09:57]}; \texttt{"classic double speak, Newspeak"} \texttt{[13:21--13:33]} \newline  \textbf{Skeptical distancing:} \texttt{"supposedly more policemen injured"} \texttt{[06:49]} & Repetition (\texttt{lie upon lie}), visceral somatic idioms (\texttt{made my blood boil}), and literary allusions (\textit{1984}, \textit{Jallianwala Bagh}, Churchill/JFK/MLK speeches at \texttt{[24:35--24:43]}) intensify moral outrage, while \texttt{"supposedly"} casts doubt on police casualty reports. \\
\bottomrule
\end{longtable}

\subsection*{Short summary (Event 2)}
\begin{itemize}[leftmargin=1.5em, itemsep=3pt, topsep=4pt]
  \item \textbf{What interpretation of the event does this video encourage?} That the 20 July lathi/tear-gas/pellet-gun crackdown on teenagers and 20-somethings marching to Parliament---after 25 days of total ministerial indifference (\texttt{[20:35--20:54]}) and Sonam Wangchuk's forced removal (\texttt{[21:14--21:23]})---is an Orwellian betrayal of Indian democracy (\texttt{"living in 1984"}, \texttt{"not the India I was taught to love"}, \texttt{[14:02]}, \texttt{[26:32]}) driven by \texttt{"arrogant autocrats"} (\texttt{[23:47]}) who fear educated youth (\texttt{[18:23--20:25]}).
  \item \textbf{Whose perspective is foregrounded, challenged, or missing?} \textit{Foregrounded:} Dissenting liberal-constitutionalist intelligentsia (Karan Thapar \& Naseeruddin Shah) standing in moral solidarity with injured NEET protesters (\texttt{Sheikh Irshad Mansuri}, \texttt{[10:19]}) and jailed student activists (\texttt{CAA} parallel via Kapil Sibal, \texttt{[26:41--27:37]}). \textit{Challenged:} PM Modi's \texttt{"Pradhan Sevak"} / \texttt{"non-biological"} persona (\texttt{[00:28]}, \texttt{[23:47--24:54]}), Dharmendra Pradhan's X post promising \texttt{"answers, reforms and accountability"} (\texttt{[12:50--13:51]}), Delhi Police casualty claims (\texttt{[06:49--07:07]}), and silent Bollywood/elite peers (\texttt{[29:18--30:00]}). \textit{Missing:} Direct voices of CJP students on the ground, government/NTA spokespersons, and injured police officers.
  \item \textbf{What evidence supports your interpretation? What complicates it?} \textit{Supports:} Named \textit{Hindu} report on 25-year-old Sheikh Irshad Mansuri's neck pellet surgery (\texttt{[10:07--10:38]}), \textit{Hindustan Times} injury reports (\texttt{[06:33]}), historical cabinet resignation precedents (\texttt{Lal Bahadur Shastri 1956}, \texttt{Madhavrao Scindia 1992}, \texttt{Nitish Kumar 1999} cited at \texttt{[22:11--22:31]}), Dharmendra Pradhan's verbatim X post (\texttt{[13:00]}), and Kapil Sibal's legal analysis of unnamed FIRs (\texttt{[26:41--27:18]}). \textit{Complicates:} Shah's confident assertion that \texttt{"no way is he [Pradhan] going to resign"} (\texttt{[12:23--12:33]}) was disproven just 72 hours later; several anti-Modi personal attacks (\texttt{"patently fake degree"}, \texttt{"uneducated himself"}, marital/family remarks at \texttt{[18:23--19:01]}) drift into ad hominem polemic far beyond NEET; and Karan Thapar never challenges his guest once, producing an echo-chamber interview rather than adversarial questioning.
\end{itemize}
\vspace{4pt}\noindent\rule{\textwidth}{0.4pt}\vspace{4pt}

\section*{Common form: Blank excerpt record}
\textit{Copy or print this page whenever you find a passage worth analysing. Keep the exact wording and enough context to interpret it.}


\subsection*{Excerpt Record 1 (Event 1 --- Connecting PMO \& Nripendra Misra to Ayodhya Land/Donation Scandal)}
\begin{itemize}[leftmargin=1.5em, itemsep=3pt, topsep=4pt]
  \item \textbf{Outlet:} The Wire | \textbf{Date:} \texttt{2026-07-01}
  \item \textbf{Video title:} \textit{Ram Temple Donations Loot: Biggest Credibility Test For Modi-RSS}
  \item \textbf{Timestamp(s):} \texttt{[11:48--13:42]} | \textbf{Analyst:} Baani Solanki | \textbf{Language:} English (Roman script)
  \item \textbf{Exact transcript excerpt:}
    \begin{quote}\small\ttfamily [11:48] [Nilanjan Mukhopadhyay:] The government of India also decided to place the gentleman who had been the principal secretary of Mr. Modi in the years between 2014 and 2019, Mr. Nripendra [11:58] Misra, to personally oversee the construction of the Ram temple... [12:17] He was very much a part of the core team of Mr. Modi. Besides that, as the person who was [12:27] supervising the construction, he would have virtually gone through and decided, played a very decisive role in deciding where the material is going to be [12:36] bought, which land is going to be bought from whom, at what price... [12:55] artificial inflation of money so that people could personally benefit on it, Mr. Misra would have been in the know of [13:04] that. And I find it very hard to believe that if he was in the know of something, the others who had actually played a decisive role in appointing him, [13:15] you know, did not know what all was happening.\end{quote}
  \item \textbf{Feature(s):} $\boxtimes$ Naming $\square$ Presupposition $\boxtimes$ Agency $\square$ Antagonisation $\boxtimes$ Intensification / softening
  \item \textbf{Who is speaking? Is this the outlet's wording or a quotation?} Senior journalist and author Nilanjan Mukhopadhyay, prompted and affirmed by The Wire host M.K. Venu (\texttt{[13:24--13:42]}).
  \item \textbf{What does the wording imply, foreground, soften, or leave unstated?} Uses deductive modal chains (\texttt{"would have virtually gone through"}, \texttt{"would have been in the know"}, \texttt{"find it very hard to believe... did not know"}) to trace institutional agency from Ayodhya land/material procurement directly into \texttt{"the core team of Mr. Modi"} at the PMO---softened grammatically as logical inference (\texttt{would have}) rather than direct documentary proof.
  \item \textbf{What context or evidence supports or complicates this interpretation?} Supported by Nripendra Misra's undisputed status as ex-Principal Secretary to PM Modi (\texttt{2014--2019}) appointed to chair the Construction Committee (\texttt{[11:48--11:58]}). Complicated by \textbf{Aaj Tak} (\texttt{B4L1pz0o4dc}, \texttt{[04:35--06:05]}, \texttt{[10:57--11:24]}), where Misra himself publicly blew the whistle on counting-room informality (\texttt{"dacoity"}) and Anjana Om Kashyap argued he oversaw \textit{construction}, not \textit{Trust treasury/counting}.
  \item \textbf{How does another outlet describe the same event or actor?} \textbf{Aaj Tak} (\texttt{B4L1pz0o4dc}, \texttt{[10:57--11:24]}) explicitly exonerates Nripendra Misra (\texttt{"nirman alag hai, prabandhan alag hai"}), whereas \textbf{The Wire} and \textbf{Dhruv Rathee} (\texttt{oeH0fOSJuX8}, \texttt{[04:28--04:46]}) treat Misra's appointment as proof that the Trust was run from Delhi/PMO.
  \item \textbf{My interpretation, with reasons:} This passage exemplifies The Wire's analytical register: instead of shouting \texttt{"chor"} (like TV panelists), Mukhopadhyay uses bureaucratic command-chain logic to pierce PMO plausible deniability.
\end{itemize}

\subsection*{Excerpt Record 2 (Event 1 --- Sangh Parivar "Shadow Organizations" \& Strategic Ambiguity)}
\begin{itemize}[leftmargin=1.5em, itemsep=3pt, topsep=4pt]
  \item \textbf{Outlet:} The Wire | \textbf{Date:} \texttt{2026-07-01}
  \item \textbf{Video title:} \textit{Ram Temple Donations Loot: Biggest Credibility Test For Modi-RSS}
  \item \textbf{Timestamp(s):} \texttt{[17:24--18:40]} \& \texttt{[24:36--25:01]} | \textbf{Analyst:} Baani Solanki | \textbf{Language:} English (Roman script)
  \item \textbf{Exact transcript excerpt:}
    \begin{quote}\small\ttfamily [17:24] [M.K. Venu:] RSS always believed in an unstructured existence... [17:43] whatever money comes to the Sangh, RSS's Guru Dakshina, therefore it is non-taxable and it can be [17:51] spent in whatever way they want to, and they need not keep any accounts. And they need not even register themselves... [18:30] would you say that this is a result of this history where the RSS or the Sangh Parivar, they refuse to be accountable?\end{quote}
    \begin{quote}\small\ttfamily [24:36] [Nilanjan Mukhopadhyay:] Setting up structures was not done. It [24:44] was probably not done because they wanted that it should not be a leak-proof system. The moment you structure something, then it becomes [24:53] leak-proof. So you keep a certain amount of ambiguity, a certain amount of fluidity.\end{quote}
  \item \textbf{Feature(s):} $\boxtimes$ Naming $\boxtimes$ Presupposition $\boxtimes$ Agency $\boxtimes$ Antagonisation $\boxtimes$ Intensification / softening
  \item \textbf{Who is speaking? Is this the outlet's wording or a quotation?} Host M.K. Venu (\texttt{[17:24--18:40]}) and guest Nilanjan Mukhopadhyay (\texttt{[24:36--25:01]}).
  \item \textbf{What does the wording imply, foreground, soften, or leave unstated?} Reframes administrative "lack of SOPs" from accidental incompetence into \textbf{purposeful institutional design} (\texttt{"they wanted that it should not be a leak-proof system"}, \texttt{"keep a certain amount of ambiguity"}), rooted in a 40-year history of unregistered Sangh bodies (\texttt{Guru Dakshina}, \texttt{Mukti Yagna Samiti}, \texttt{[17:43]}, \texttt{[20:46--21:20]}). Note the careful epistemic hedge \texttt{"probably"} (\texttt{[24:44]}) softening an otherwise explosive claim of intentionality.
  \item \textbf{What context or evidence supports or complicates this interpretation?} Supported by the 1990 \texttt{Rs.\,700 crore} VHP tax-notice episode (\texttt{[15:33--17:15]}) and the November 2020 private audit warning that went unimplemented for 6 years (cited in Dhruv Rathee \texttt{oeH0fOSJuX8} \texttt{[10:12--10:32]} and Republic TV panelist Ashish Chaturvedi \texttt{0H-bruoEkaY} \texttt{[27:34--27:53]}).
  \item \textbf{How does another outlet describe the same event or actor?} \textbf{Aaj Tak}'s VHP spokesperson (\texttt{B4L1pz0o4dc}, \texttt{[34:21]}) calls the exact same lack of structure harmless \texttt{"kuch avyavastha"} [some disarray], while \textbf{Republic TV} (\texttt{0H-bruoEkaY}, \texttt{[26:56]}) attributes it merely to 79-year-old Champat Rai being \texttt{"not skilled administratively"}.
  \item \textbf{My interpretation, with reasons:} Comparing \texttt{"kuch avyavastha"} (Aaj Tak/VHP) with \texttt{"they wanted that it should not be a leak-proof system"} (The Wire) shows how two outlets observe the \textit{exact same absence of accounting controls} and assign opposite moral agency.
\end{itemize}

\subsection*{Excerpt Record 3 (Event 1 --- "Ayodhya the Hindu Vatican" \& Hollowed-Out Devotee Faith)}
\begin{itemize}[leftmargin=1.5em, itemsep=3pt, topsep=4pt]
  \item \textbf{Outlet:} The Wire | \textbf{Date:} \texttt{2026-07-01}
  \item \textbf{Video title:} \textit{Ram Temple Donations Loot: Biggest Credibility Test For Modi-RSS}
  \item \textbf{Timestamp(s):} \texttt{[35:04--37:08]} | \textbf{Analyst:} Baani Solanki | \textbf{Language:} English (Roman script)
  \item \textbf{Exact transcript excerpt:}
    \begin{quote}\small\ttfamily [35:04] [Nilanjan Mukhopadhyay:] Ayodhya in a way, Ayodhya the [35:13] temple, Ayodhya the Hindu Vatican is actually the Sangh Parivar's creation. So, within the Hindu Vatican, it does [35:22] not give any space to these established, you know, our temporal [spiritual] heads... [36:01] the moment people lose their rationality, they become bhakts. [36:11] We really do not know to what extent people have become bhakts. That they are willing to [36:20] overlook anything that happens... [36:28] they put their body and soul into this movement. You remember the Ram Shila Yatras of 1989... [36:59] And now you find that the darshan has been hollowed out, that it has been stolen.\end{quote}
  \item \textbf{Feature(s):} $\boxtimes$ Naming $\boxtimes$ Presupposition $\square$ Agency $\boxtimes$ Antagonisation $\boxtimes$ Intensification / softening
  \item \textbf{Who is speaking? Is this the outlet's wording or a quotation?} Nilanjan Mukhopadhyay speaking in dialogue with M.K. Venu.
  \item \textbf{What does the wording imply, foreground, soften, or leave unstated?} Coins the striking metaphor \texttt{"Ayodhya the Hindu Vatican"} (\texttt{[35:13]}) to argue that the Sangh Parivar centralized a pluralistic Hindu tradition (historically centred in Varanasi, Mathura, Haridwar) into a politically controlled headquarters while sidelining traditional \textit{acharyas} (\texttt{[34:07--35:22]}). Simultaneously distinguishes blind \texttt{"bhakts"} (\texttt{"lose their rationality"}, \texttt{[36:01]}) from sincere grassroots \texttt{1989} \textit{Ram Shila} devotees whose sacred \texttt{"darshan has been hollowed out"} (\texttt{[36:59]}).
  \item \textbf{What context or evidence supports or complicates this interpretation?} Supported by public criticism from traditional \textit{akhadas} (\texttt{Nirmohi Akhada} at \texttt{[29:20]}) and \textit{Shankaracharyas} (also referenced by SP's Anurag Bhadauria on Aaj Tak \texttt{B4L1pz0o4dc} \texttt{[15:17]}). Complicated by Mukhopadhyay's own candid admission (\texttt{[37:08--37:19]}) that \texttt{"We do not have any empirical data as yet"} on whether devotees will actually turn against the BJP.
  \item \textbf{How does another outlet describe the same event or actor?} \textbf{Sarthak Goswami} (\texttt{lb7gZnyZH0c}, \texttt{[06:26]}) uses a remarkably parallel Hindi formulation---\texttt{"Ram ka chanda nahin Ram ka charitra chori ho chuka hai"}---to express the same "hollowed-out sacredness" in popular Hinglish rather than academic English.
  \item \textbf{My interpretation, with reasons:} By separating pluralistic Hinduism and sincere \textit{Ram Shila} pilgrims from \texttt{"the Hindu Vatican"} of the Sangh Parivar, The Wire critiques Hindutva institutional power while empathizing with ordinary believers' sense of betrayal.
\end{itemize}

\subsection*{Excerpt Record 4 (Event 2 --- Unbadged "Army of Goons", Neck Pellets \& Student Retaliation)}
\begin{itemize}[leftmargin=1.5em, itemsep=3pt, topsep=4pt]
  \item \textbf{Outlet:} The Wire (\textit{The Interview}) | \textbf{Date:} \texttt{2026-07-22}
  \item \textbf{Video title:} \textit{I stand with the protesting students; this is not the India I love; I don't take Modi seriously}
  \item \textbf{Timestamp(s):} \texttt{[07:15--07:52]}, \texttt{[09:36--09:57]} \& \texttt{[10:07--10:38]} | \textbf{Analyst:} Baani Solanki | \textbf{Language:} English (Roman script)
  \item \textbf{Exact transcript excerpt:}
    \begin{quote}\small\ttfamily [07:15] [Naseeruddin Shah:] these people who are sadistically beating up these students who are asking [07:25] for their rights are young people themselves most of the time. A lot of them don't look like policemen... [07:34] they're not carrying name badges, they're not carrying ranks, they're not even in proper uniform. Some of them are even admitting that they are private [07:42] privately hired. So to let loose an army of goons on our future...\end{quote}
    \begin{quote}\small\ttfamily [09:36] And what I watched this morning is the students retaliating. It's not surprising. I mean, even a rat fights back when it's [09:46] pushed against the wall. And these kids are justified in what they're asking. Their anger is absolutely valid...\end{quote}
    \begin{quote}\small\ttfamily [10:07] [Karan Thapar:] today's Hindu reports that pellet guns and shock batons were used on some of the students. [10:19] It reports that Sheikh Irshad Mansuri, who's only 25 years old, had to undergo surgery for pellets lodged in the upper part of his body and his neck to be removed.\end{quote}
  \item \textbf{Feature(s):} $\boxtimes$ Naming $\boxtimes$ Presupposition $\boxtimes$ Agency $\boxtimes$ Antagonisation $\boxtimes$ Intensification / softening
  \item \textbf{Who is speaking? Is this the outlet's wording or a quotation?} Guest Naseeruddin Shah (\texttt{[07:15--09:57]}) and The Wire interviewer Karan Thapar citing \textit{The Hindu} (\texttt{[10:07--10:38]}).
  \item \textbf{What does the wording imply, foreground, soften, or leave unstated?} Combines Shah's moral-emotional intensification (\texttt{"sadistically beating"}, \texttt{"an army of goons on our future"}) with Karan Thapar's hard newspaper citation (\textit{The Hindu} on 25-year-old \texttt{Sheikh Irshad Mansuri}'s neck pellet surgery). Crucially, student counter-violence is explicitly softened and morally excused via \texttt{"retaliating"} and the cornered-victim idiom \texttt{"even a rat fights back when it's pushed against the wall"} (\texttt{[09:36--09:46]}).
  \item \textbf{What context or evidence supports or complicates this interpretation?} \textit{The Hindu}'s medical detail on Sheikh Irshad Mansuri provides concrete verification of pellet-gun injuries cited more generally by \textbf{Sarthak Goswami} (\texttt{EM1GeZufO6E}, \texttt{[03:39--03:47]}). Meanwhile, Shah's dismissal of police injuries at \texttt{[06:49--07:07]} (\texttt{"supposedly more policemen injured... something like 600"}) exaggerates/garbles the official \texttt{110--118} figure documented with X-rayed fractures on \textbf{Aaj Tak} (\texttt{JXLIs0WtJAM}, \texttt{[18:09]}, \texttt{[21:12--23:23]}).
  \item \textbf{How does another outlet describe the same event or actor?} Notice the exact lexical mirror image with \textbf{Republic TV} (\texttt{HTvb6E\_1Qdc}, \texttt{[04:19]}): Arnab calls protesters \texttt{"the goons of Abhijit Dipke"} attacking uniformed officers, whereas Naseeruddin Shah calls unbadged state-aligned assailants \texttt{"an army of goons"} let loose on \texttt{"our future"}!
  \item \textbf{My interpretation, with reasons:} Pairing Karan Thapar's print-journalism verification (\textit{The Hindu}) with Shah's visceral moral outrage gives The Wire both factual anchoring and emotional force, while clearly asymmetrically excusing student \texttt{"retaliating"}.
\end{itemize}

\subsection*{Excerpt Record 5 (Event 2 --- Deconstructing Dharmendra Pradhan's X Post as Orwellian "Double Speak")}
\begin{itemize}[leftmargin=1.5em, itemsep=3pt, topsep=4pt]
  \item \textbf{Outlet:} The Wire (\textit{The Interview}) | \textbf{Date:} \texttt{2026-07-22}
  \item \textbf{Video title:} \textit{I stand with the protesting students; this is not the India I love; I don't take Modi seriously}
  \item \textbf{Timestamp(s):} \texttt{[12:50--14:20]} | \textbf{Analyst:} Baani Solanki | \textbf{Language:} English (Roman script)
  \item \textbf{Exact transcript excerpt:}
    \begin{quote}\small\ttfamily [12:50] [Karan Thapar:] Yesterday the education minister posted on X. This is one of the [13:00] things he said. Speaking to the students, he wrote: "We owe them answers, reforms and accountability. This is exactly what our government [13:10] remains committed to delivering." Accountable is the one thing the minister refuses to be and his silence suggests he refuses to answer their [13:21] questions. How do you view his statement?\end{quote}
    \begin{quote}\small\ttfamily [Naseeruddin Shah:] This is classic double speak. What it was called in 1984, the book, [13:33] classic double speak, Newspeak... He's sticking to the vocabulary which he hopes will mollify some... [14:02] we are living in 1984 where there's surveillance on everybody. 2 minutes of hate has been transformed [14:11] into 24 hours of hate. And it's just completely turning [14:20] into a dystopian scenario here.\end{quote}
  \item \textbf{Feature(s):} $\boxtimes$ Naming $\boxtimes$ Presupposition $\square$ Agency $\boxtimes$ Antagonisation $\boxtimes$ Intensification / softening
  \item \textbf{Who is speaking? Is this the outlet's wording or a quotation?} Karan Thapar quoting Dharmendra Pradhan's official X post (\texttt{[13:00--13:10]}) and adding his own editorial commentary (\texttt{[13:10--13:21]}), followed by Naseeruddin Shah's analysis (\texttt{[13:21--14:20]}).
  \item \textbf{What does the wording imply, foreground, soften, or leave unstated?} Notice how Karan Thapar pre-frames the quote \textit{before} asking his question (\texttt{"Accountable is the one thing the minister refuses to be"}---leaving zero interpretive neutrality). Shah then performs explicit \textbf{critical discourse analysis} using George Orwell's \textit{1984} (\texttt{"classic double speak, Newspeak"}, \texttt{"24 hours of hate"}), arguing that state actors deploy hollow democratic buzzwords (\texttt{answers, reforms, accountability}) while practising surveillance and repression.
  \item \textbf{What context or evidence supports or complicates this interpretation?} Supported by the glaring 30-day gap during which Dharmendra Pradhan never once visited protesting students at Jantar Mantar (\texttt{[20:35--20:54]}, corroborated by citizen interviews in Mohak Mangal \texttt{v-RikSqb0Vw} \texttt{[05:31--05:38]}). Complicated by the fact that 3 days after this interview (\texttt{25 July}), Pradhan \textit{did} resign (\texttt{EM1GeZufO6E} \texttt{[00:08]}), showing democratic pressure still functioned beyond a total \textit{1984} dystopia.
  \item \textbf{How does another outlet describe the same event or actor?} \textbf{Mohak Mangal} (\texttt{v-RikSqb0Vw}, \texttt{[00:00--01:03]}) critiques Dharmendra Pradhan not through Western literary theory (\textit{Orwell's 1984}), but through Indian biographical irony (his \texttt{1997} Odisha ABVP lathi-charge injuries)---appealing to a much broader Hindi/Hinglish youth audience.
  \item \textbf{My interpretation, with reasons:} This excerpt showcases The Wire's cultural register---using cosmopolitan English literary touchstones (\textit{Orwell's 1984}, \texttt{Newspeak}) and pre-loaded interviewer prompts to indict official government communication.
\end{itemize}

\subsection*{Excerpt Record 6 (Event 2 --- Citizens' Right to Parliament \& Kapil Sibal's UAPA Warning)}
\begin{itemize}[leftmargin=1.5em, itemsep=3pt, topsep=4pt]
  \item \textbf{Outlet:} The Wire (\textit{The Interview}) | \textbf{Date:} \texttt{2026-07-22}
  \item \textbf{Video title:} \textit{I stand with the protesting students; this is not the India I love; I don't take Modi seriously}
  \item \textbf{Timestamp(s):} \texttt{[25:02--27:37]} | \textbf{Analyst:} Baani Solanki | \textbf{Language:} English (Roman script)
  \item \textbf{Exact transcript excerpt:}
    \begin{quote}\small\ttfamily [25:02] [Karan Thapar:] Do peaceful nonviolent Indian protesters protesting in a just cause have the [25:13] right to walk to parliament and demand the attention of their MPs and government? The government's answer is clearly they don't. In which case, how [25:23] meaningful is our democracy?\end{quote}
    \begin{quote}\small\ttfamily [Naseeruddin Shah:] That question answers itself. Karan, this parliament is run with our [25:33] taxpayer money. These people are our representatives. That house of parliament is ours. It belongs to the kids. It belongs to the [25:42] future. They have an absolute right to enter it... [26:32] This is not the India I was taught to love.\end{quote}
    \begin{quote}\small\ttfamily [26:41] [Karan Thapar:] The police have already filed FIRs against unnamed persons. Kapil Sibal fears that in due course [26:58] they will find names and possibly charge people under UAPA or section 152 of the old penal code [27:07] which is now section 195 of BNS... that's exactly what happened 6 years ago when [27:18] the CAA protests occurred. And those young people from the CAA days are still rotting [ASR: pushing] in jail.\end{quote}
  \item \textbf{Feature(s):} $\boxtimes$ Naming $\boxtimes$ Presupposition $\boxtimes$ Agency $\boxtimes$ Antagonisation $\boxtimes$ Intensification / softening
  \item \textbf{Who is speaking? Is this the outlet's wording or a quotation?} Karan Thapar (\texttt{[25:02--25:23]}, \texttt{[26:41--27:27]}) and Naseeruddin Shah (\texttt{[25:23--26:32]}).
  \item \textbf{What does the wording imply, foreground, soften, or leave unstated?} Reclaims democratic ownership of physical state space (\texttt{"That house of parliament is ours. It belongs to the kids"}, \texttt{[25:33--25:42]}) against security-state exclusion, and links the 5 open-ended Delhi Police FIRs (reported positively on Aaj Tak \texttt{JXLIs0WtJAM} \texttt{[19:25--20:22]}) to a chilling pattern of post-protest criminalization (\texttt{UAPA} / \texttt{BNS 195} pre-trial detention since the 2020 CAA protests).
  \item \textbf{What context or evidence supports or complicates this interpretation?} Supported by the actual filing of 5 FIRs with unnamed accused (\texttt{Aaj Tak} \texttt{[19:25--20:22]}) and by CJP making \textit{withdrawal/non-filing of police cases} one of its core settlement conditions before ending the protest on 25 July (\texttt{Republic TV} \texttt{HTvb6E\_1Qdc} \texttt{[00:07--00:15]}; \texttt{Sarthak Goswami} \texttt{EM1GeZufO6E} \texttt{[08:25]})! Complicated legally by the fact that even in constitutional democracies, unrestricted mass entry into an active Parliament building (\texttt{"an absolute right to enter it"}, \texttt{[25:42]}) is subject to reasonable security regulation.
  \item \textbf{How does another outlet describe the same event or actor?} \textbf{Aaj Tak} (\texttt{JXLIs0WtJAM}, \texttt{[19:25--20:41]}) presents the exact same 5 FIRs and Parliament cordons neutrally as routine \texttt{"ligal ekshan"} around \texttt{"vaital instoleshanans"}, whereas \textbf{The Wire} exposes how open-ended FIRs can become multi-year \texttt{UAPA} traps for students.
  \item \textbf{My interpretation, with reasons:} Highlighting the danger of unnamed FIRs proved prescient---explaining why CJP refused to vacate Jantar Mantar even after Pradhan resigned until the government explicitly guaranteed no cases would be pursued (\texttt{HTvb6E\_1Qdc}, \texttt{[00:07--00:15]}).
\end{itemize}
\vspace{4pt}\noindent\rule{\textwidth}{0.4pt}\vspace{4pt}

\section*{Common form: End-of-event summary}

\subsection*{End-of-Event Summary --- Event 1 (Ram Mandir Donation Controversy)}
\begin{itemize}[leftmargin=1.5em, itemsep=3pt, topsep=4pt]
  \item \textbf{The main way this outlet framed the event was:} As a historic moral and institutional unraveling of the Sangh Parivar's core 40-year project (\texttt{"Biggest Credibility Test For Modi-RSS"}, \texttt{"biggest challenge... after the assassination of Mahatma Gandhi"}, \texttt{[08:23--08:31]}), arguing that massive donation \texttt{"dacoity"} and \texttt{5x--10x} land-price \textit{hawala} inflation stemmed directly from RSS's unaccountable \texttt{"shadow organization"} culture (\texttt{[17:24--22:32]}), intentional avoidance of \texttt{"leak-proof"} structures (\texttt{[24:44--25:01]}), and PMO-linked oversight under Nripendra Misra (\texttt{[11:48--14:55]}).
  \item \textbf{The strongest linguistic evidence was (quote + timestamp):} \texttt{"Setting up structures was not done. It was probably not done because they wanted that it should not be a leak-proof system... Ayodhya the Hindu Vatican is actually the Sangh Parivar's creation... And now you find that the darshan has been hollowed out, that it has been stolen."} (\texttt{[24:36--24:53]}, \texttt{[35:13]} \& \texttt{[36:59]}).
  \item \textbf{The shared feature(s) most clearly illustrated were:} \textbf{Naming} (\texttt{shadow organizations}, \texttt{corporate bureaucratic structure}, \texttt{political loyalists}, \texttt{Ayodhya the Hindu Vatican}), \textbf{Agency} (attributing lack of accounting controls to deliberate Sangh design rather than clerical oversight at \texttt{[24:44--25:01]}), and \textbf{Epistemic Hedging vs Historical Intensification} (\texttt{potentially}, \texttt{no empirical data} paired with \texttt{100, maybe 1,000 questions} and \texttt{1948} parallels).
  \item \textbf{An alternative interpretation of the same evidence is:} That transitioning from an agitational movement to managing millions of pilgrims created genuine early administrative bottlenecks rather than intentional \texttt{"non-leak-proof"} conspiracy, and that Nripendra Misra's public admission of \texttt{"dacoity"} alongside 8 SIT arrests and senior resignations reflected corrective institutional action rather than PMO complicity.
  \item \textbf{What I cannot conclude from this transcript alone:} Who the ultimate financial beneficiaries of the inflated land transactions were (\texttt{[04:48]}), whether Nripendra Misra actually briefed PM Modi on donation irregularities, or whether voter outrage will translate into electoral losses in UP (\texttt{[37:08--37:28]}).
\end{itemize}

\subsection*{End-of-Event Summary --- Event 2 (Students' Protest Regarding NEET Paper Leak)}
\begin{itemize}[leftmargin=1.5em, itemsep=3pt, topsep=4pt]
  \item \textbf{The main way this outlet framed the event was:} As an agonizing democratic and moral crisis (\texttt{"this is not the India I was taught to love"}, \texttt{"living in 1984"}, \texttt{[14:02]}, \texttt{[26:32]}) in which \texttt{"arrogant autocrats"} (\texttt{[23:47]}) ignored wronged students for 25 days, offered Orwellian \texttt{"double speak"} (\texttt{[13:21]}), and unleashed tear gas, neck-lodging pellet guns (\texttt{Sheikh Irshad Mansuri}, \texttt{[10:19]}), and unbadged \texttt{"lumpens"} / \texttt{"an army of goons"} (\texttt{[07:42]}, \texttt{[11:26]}) on peaceful teenagers rightfully marching toward citizens' Parliament.
  \item \textbf{The strongest linguistic evidence was (quote + timestamp):} \texttt{"to let loose an army of goons on our future... This is classic double speak. What it was called in 1984... That house of parliament is ours. It belongs to the kids. It belongs to the future."} (\texttt{[07:42]}, \texttt{[13:21--13:33]} \& \texttt{[25:33--25:42]}).
  \item \textbf{The shared feature(s) most clearly illustrated were:} \textbf{Presupposition} (Karan Thapar's heavily loaded negative-interrogative and declarative prompts at \texttt{[06:41]}, \texttt{[13:10]}, \texttt{[22:40]}, \texttt{[23:47]}), \textbf{Naming} (\texttt{army of goons} and \texttt{lumpens} for state-aligned attackers vs \texttt{these kids} / \texttt{these bachas} / \texttt{our future} for protesters), and \textbf{Agency} (framing student clashes strictly as justified cornered self-defense---\texttt{"even a rat fights back when it's pushed against the wall"}, \texttt{[09:36--09:46]}).
  \item \textbf{An alternative interpretation of the same evidence is:} That security cordons prevented a chaotic physical breach of Parliament during session, that over \texttt{110} uniformed police officers suffered genuine fractures/injuries from crowd stone-pelting (dismissed skeptically as \texttt{"supposedly"} by Shah at \texttt{[06:49]}), and that Indian democratic accountability was far less closed than \textit{1984}---proven when Dharmendra Pradhan resigned just 3 days after this interview.
  \item \textbf{What I cannot conclude from this transcript alone:} What protesting students on the ground prioritized in negotiations (since no student was interviewed), whether the unbadged men were formally police personnel or privately paid recruits (\texttt{[07:34--07:42]}), or why the government reversed course to accept CJP's demands on 25 July.
\end{itemize}
\textit{Research reminder: interpret wording in context, distinguish allegations from established facts, and support claims with exact excerpts. Mixed or contradictory evidence is worth recording.}


\end{document}
